\begin{textsample}{POS Dim 4 – ap – Score 8.00 – ap/ap\_inf\_e\_ef\_8.txt}  \label{ex:f4_pos_236}
2.4.
Modalidades de Ensino – Educação \textbf{Indígena} A Constituição de 1988 inaugurou um novo momento nas relações entre \textbf{Estado} brasileiro e os povos \textbf{indígenas} ao reconhecer e valorizar a sociodiversidade \textbf{indígena}.
Para intensificar essa política, reforçando as disposições da Constituição, a portaria interministerial do Ministério da Justiça n.o 559 de 16/04/91, passou a tratar as seguintes questões : garantia da oferta da Educação Escolar \textbf{indígena} de qualidade, \textbf{laica} e diferenciada ; ensino \textbf{bilíngue} ; criação de órgãos normativos para o acompanhamento e desenvolvimento da educação \textbf{indígena} ; recursos financeiros ; formação e capacitação de professores ; reconhecimento das instituições escolares ; garantias de continuação dos estudos em escolas comuns quando este não for oferecido nas escolas \textbf{indígenas} ; garantia de acesso ao material didático específico ; isonomia salarial entre professores \textbf{índios} e não \textbf{índios} ; e determinação da revisão da imagem do \textbf{índio}, historicamente distorcida a ser divulgada nas redes de ensino.
No âmbito da Educação Escolar ficaram consagrados os princípios da interculturalidade, valorização das línguas maternas, da atenção as diferentes realidades sociolinguísticas e as práticas curriculares e pedagógicas.
A escola desta modalidade tem uma realidade singular, inscrita em terras e culturas \textbf{indígenas}.
Requer, portanto, pedagogia própria em respeito à especificidade étnico-cultural de cada povo ou comunidade e formação específica de seu quadro docente, observados os princípios constitucionais, a Base Nacional Comum e os princípios que orientam a Educação Básica brasileira ( artigos 5o, 9o, 10, 11 e inciso VIII do artigo 4o da LDB ).
Na estruturação e no funcionamento das escolas \textbf{indígenas} é reconhecida sua condição de escolas com normas e ordenamento jurídico próprios, com ensino intercultural e \textbf{bilíngue}, visando à valorização plena das culturas dos povos \textbf{indígenas} e à afirmação e manutenção de sua diversidade étnica.
A Educação Escolar \textbf{Indígena} normatizada pela Resolução CNE/CEB no 5, de 22 de junho de 2012, define Diretrizes Curriculares Nacionais para a Educação Escolar \textbf{Indígena} na Educação Básica.
Conforme o art. 2o, das Diretrizes Curriculares Nacionais para a Educação Escolar \textbf{Indígena} na Educação Básica, têm por objetivos : I.
Orientar as escolas \textbf{indígenas} de educação básica e os sistemas de ensino da União, dos Estados, do Distrito Federal e dos Municípios na elaboração, desenvolvimento e avaliação de seus projetos educativos ; II.
Orientar os processos de construção de instrumentos normativos dos sistemas de ensino visando tornar a Educação Escolar \textbf{Indígena} projeto orgânico, articulado e sequenciado de Educação Básica entre suas diferentes etapas e modalidades, sendo garantidas as especificidades dos processos educativos \textbf{indígenas} ; III.
Assegurar que os princípios da especificidade, do bilinguismo e multilinguismo, da organização comunitária e da interculturalidade fundamentem os projetos educativos das comunidades \textbf{indígenas}, valorizando suas línguas e conhecimentos tradicionais ; IV.
Assegurar que o modelo de organização e gestão das escolas \textbf{indígenas} leve em consideração as práticas socioculturais e econômicas das respectivas comunidades, bem como suas formas de produção de conhecimento, processos próprios de ensino e de aprendizagem e projetos societários ; V.
Fortalecer o regime de colaboração entre os sistemas de ensino da União, dos Estados, do Distrito Federal e dos Municípios, fornecendo diretrizes para a organização da Educação Escolar \textbf{Indígena} na Educação Básica, no âmbito dos \textbf{territórios} etnoeducacionais ; VI.
Normatizar dispositivos constantes na Convenção 169, da Organização Internacional do Trabalho, ratificada no Brasil, por meio do Decreto Legislativo no 143/2003, no que se refere à educação e meios de comunicação, bem como os mecanismos de consulta livre, prévia e informada ; VII.
Orientar os sistemas de ensino da União, dos Estados, do Distrito Federal e dos Municípios a incluir, tanto nos processos de formação de professores \textbf{indígenas}, quanto no funcionamento regular da Educação Escolar \textbf{Indígena}, a colaboração e atuação de especialistas em saberes tradicionais, como os tocadores de instrumentos musicais, contadores de narrativas míticas, pajés e xamãs, rezadores, raizeiros, parteiras, organizadores de rituais, conselheiros e outras funções próprias e necessárias ao bem viver dos povos \textbf{indígenas} ; VIII.
Zelar para que o direito à educação escolar diferenciada seja garantido às comunidades \textbf{indígenas} com qualidade social e pertinência pedagógica, cultural, linguística, ambiental e territorial, respeitando as lógicas, saberes e perspectivas dos próprios povos \textbf{indígenas}.
Além do reconhecimento do direito dos \textbf{índios} de manterem sua \textbf{identidade} cultural, cabe ao \textbf{Estado} proteger as manifestações culturais \textbf{indígenas} e o uso de suas línguas maternas.
Esses dispositivos abriram a possibilidade para que as escolas \textbf{indígenas} se constituam como instrumento de valorização das línguas, dos saberes e das tradições, tornando -se espaço de diálogos entre culturas diferentes.
De acordo com o texto das Diretrizes Curriculares para a Educação \textbf{Indígena} que se reporta ao Parecer CNE/CEB no 14/99, “ reconhece que a escola \textbf{indígena} é uma experiência pedagógica peculiar e como tal deve ser tratada pelas agências governamentais promovendo as adequações institucionais e legais necessárias para garantir a implementação de uma política de governo que priorize assegurar às sociedades \textbf{indígenas} uma educação diferenciada, respeitando seu universo sócio cultural ”.
Desta forma, cumprindo com o preceito legal estabelecido pelo referido Parecer, que a escola \textbf{indígena} adquira características próprias e peculiares tais como : Específica e diferenciada : porque concebida e planejada como reflexo das aspirações particulares de cada povo \textbf{indígena} e com autonomia com relação a determinados aspectos que regem o funcionamento e orientação da escola não \textbf{indígena} ; Intercultural : porque deve reconhecer e manter a diversidade cultural e linguística ; promover uma situação de comunicação entre experiências socioculturais, linguísticos e históricos diferentes, não considerando uma cultura superior a outra, estimular o entendimento e o respeito entre seres humanos de \textbf{identidades} étnicas diferentes, ainda que se reconheça que tais relações vêm ocorrendo historicamente em contextos de desigualdade social e política.
Bilíngue/multilíngue : porque as tradições culturais, os conhecimentos acumulados, a educação das gerações mais novas, as crenças, o pensamento e a prática \textbf{religiosa}, enfim as reproduções socioculturais das sociedades \textbf{indígenas} são na maioria dos casos manifestadas através do uso de mais de uma língua, [... ] constituindo assim um quadro de bilínguismo. ( p.24 e 25 RCNEI. 2002 ) Assim é a escola \textbf{indígena}, o direito ao exercício e ao desenvolvimento das próprias culturas, aos seus conhecimentos tradicionais e ao direito de ofertar às suas crianças a alfabetização na própria língua de acordo com seus processos próprios de ensino e aprendizagem.
Os currículos das escolas \textbf{indígenas}, construídos por seus professores em articulação com as comunidades \textbf{indígenas} deverão ser aprovados pelos respectivos órgãos normativos dos sistemas de ensino.
Desta forma deverão ser acrescidos ao currículo os seguintes componentes curriculares : Cultura \textbf{Indígena} e Língua Materna, cujos conteúdos serão elaborados pelos professores \textbf{indígenas} de cada etnia.
O artigo 15, da Resolução CNE/CEB no 05 de 22/06/12, determina que “ O currículo das escolas \textbf{indígenas}, ligado às concepções e práticas que definem o papel sociocultural da escola, diz respeito aos modos de organização dos tempos e espaços da escola, de suas atividades pedagógicas, das relações sociais tecidas no cotidiano escolar, das interações do ambiente educacional com a sociedade, das relações de poder presentes no fazer educativo e nas formas de conceber e construir conhecimentos escolares, constituindo parte importante dos processos sociopolíticos e culturais de construção de \textbf{identidades}. ”

% matched lemmas: bilíngue, estado, identidade, indígena, laico, religioso, território, índio
\end{textsample}
